%-------------------------
% Resume in Latex
% Author : Chenshuhao Qin
% v2026.10 — Infra/data-platform pivot of v2026.08: Big Data & Lakehouse skills line,
%           Spark/Flink/Iceberg/Kafka + ADL Foundry + K8s Operator bullets at Alibaba,
%           trace-analysis framing, CUDA, availability line in header
%------------------------

\documentclass[letterpaper,10pt]{article}

\usepackage{latexsym}
\usepackage[empty]{fullpage}
\usepackage{titlesec}
\usepackage{marvosym}
\usepackage[usenames,dvipsnames]{color}
\usepackage{verbatim}
\usepackage{enumitem}
\usepackage[hidelinks]{hyperref}
\usepackage{fancyhdr}
\usepackage[english]{babel}
\usepackage{tabularx}
\input{glyphtounicode}

% Font options
\usepackage{charter}

\pagestyle{fancy}
\fancyhf{}
\fancyfoot{}
\renewcommand{\headrulewidth}{0pt}
\renewcommand{\footrulewidth}{0pt}

% Optimizing Margins
\addtolength{\oddsidemargin}{-0.55in}
\addtolength{\evensidemargin}{-0.55in}
\addtolength{\textwidth}{1.1in}
\addtolength{\topmargin}{-.7in}
\addtolength{\textheight}{1.35in}

\urlstyle{same}

% Uniform, overlap-safe leading
\linespread{0.98}

\raggedbottom
\raggedright
\setlength{\tabcolsep}{0in}

% Section Formatting
\titleformat{\section}{
  \vspace{-8pt}\scshape\raggedright\large
}{}{0em}{}[\color{black}\titlerule \vspace{-4pt}]

% Custom Commands
\newcommand{\resumeItem}[1]{
  \item\small{
    {#1}
  }
}

\newcommand{\resumeSubheading}[4]{
  \item
    \begin{tabular*}{0.97\textwidth}[t]{l@{\extracolsep{\fill}}r}
      \textbf{#1} & #2 \\
      \textit{\small#3} & \textit{\small #4} \\
    \end{tabular*}\vspace{-1pt}
}

% Paragraph-style entry text wrapped in a bullet-less itemize so LaTeX
% guarantees safe baseline separation from the heading above and the next entry below
\newcommand{\resumeParagraph}[1]{
  \begin{itemize}[leftmargin=0.15in, label={}, itemsep=0pt, topsep=2pt, parsep=0pt]
    \item{\small #1}
  \end{itemize}
}

\newcommand{\resumeProjectHeading}[2]{
    \item
    \begin{tabular*}{0.97\textwidth}{l@{\extracolsep{\fill}}r}
      \small#1 & #2 \\
    \end{tabular*}\vspace{-1pt}
}

\renewcommand\labelitemii{$\vcenter{\hbox{\tiny$\bullet$}}$}

\newcommand{\resumeSubHeadingListStart}{\begin{itemize}[leftmargin=0.15in, label={}, itemsep=2pt, topsep=2pt, parsep=0pt]}
\newcommand{\resumeSubHeadingListEnd}{\end{itemize}}
\newcommand{\resumeItemListStart}{\begin{itemize}[leftmargin=0.15in, itemsep=1pt, topsep=2pt, parsep=0pt]}
\newcommand{\resumeItemListEnd}{\end{itemize}\vspace{-3pt}}

%-------------------------------------------
%%%%%%  RESUME STARTS HERE  %%%%%%%%%%%%%%%%%%%%%%%%%%%%

\begin{document}

%----------HEADER----------
\begin{flushleft}
    \textbf{\Huge \scshape Chenshuhao (Cody) Qin} \\ \vspace{3pt}
    \small
    Los Angeles, CA $|$ (213) 922-0925 $|$
    \href{mailto:chenshuh@usc.edu}{\underline{chenshuh@usc.edu}} $|$
    \href{https://codyqin.github.io/CodyQin/}{\underline{codyqin.github.io/CodyQin}} $|$
    \href{https://github.com/CodyQin}{\underline{github.com/CodyQin}} $|$
    \href{https://www.linkedin.com/in/chenshuhao-qin/}{\underline{linkedin.com/in/chenshuhao-qin}} \\
    \vspace{2pt}
    \textbf{Availability}: Summer 2027 internship, May -- Aug. 2027
\end{flushleft}

%-----------EDUCATION-----------
\section{Education}
  \resumeSubHeadingListStart
    \resumeSubheading
      {University of Southern California (USC)}{Los Angeles, CA}
      {Master of Science in Computer Science}{Jan. 2026 -- Dec. 2027 (Expected)}
      \resumeItemListStart
        \resumeItem{\textbf{Coursework}: Machine Learning, Deep Learning, Natural Language Processing, Data Structures \& Algorithms, Operating Systems, Computer Architecture, Computer Networks, Predictive Modeling}
      \resumeItemListEnd
    \resumeSubheading
      {Duke University; Duke Kunshan University}{Durham, NC / Kunshan, China}
      {Bachelor of Science in Computer Science; B.S. in Computation and Design}{Aug. 2021 -- May. 2025}
      \resumeItemListStart
        \resumeItem{\textbf{Honors}: \textbf{Apple WWDC23 Swift Student Challenge Winner} (1 of 13 in China), Provincial Innovation Grant Recipient (\$3k), DKU Student Ambassador}
      \resumeItemListEnd
  \resumeSubHeadingListEnd

%-----------SKILLS-----------
\section{Skills}
 \begin{itemize}[leftmargin=0.15in, label={}, itemsep=0pt, topsep=1pt, parsep=0pt]
    \small{\item{
     \textbf{Languages}{: Python, C, Java, SQL, TypeScript/JavaScript, Swift} \\
     \textbf{Big Data \& Lakehouse}{: Apache Spark, Flink, Iceberg, Kafka, Ray, batch \& streaming pipelines, object storage (OSS/S3)} \\
     \textbf{Infrastructure \& Cloud}{: Kubernetes (incl. Operator development), Docker, Linux, Git, CI/CD, MySQL, PostgreSQL, Redis, REST APIs (FastAPI, Flask), React, AWS} \\
     \textbf{AI \& Machine Learning}{: PyTorch, TensorFlow, CUDA GPU programming, LoRA fine-tuning, vLLM, Hugging Face, Pandas, NumPy} \\
     \textbf{LLM \& Agents}{: agent harness \& context engineering, retrieval-augmented generation (RAG), LangChain, LlamaIndex, AgentScope, Milvus vector search}
    }}
 \end{itemize}

%-----------WORK EXPERIENCE-----------
\section{Work Experience}
  \resumeSubHeadingListStart

    \resumeSubheading
      {Alibaba Cloud (Top-4 Global Cloud Provider)}{Hangzhou, China}
      {AI Engineer Intern (Agent Infrastructure)}{Jun. 2026 -- Aug. 2026}
      \resumeItemListStart
        \resumeItem{Collaborated with Alibaba's \textbf{Tongyi Lab} to productize \textbf{Data-Juicer Agents} --- an open-source \textbf{agent harness} --- into natural-language multimodal data pipelines (cleaning, filtering, deduplication) for the broader \textbf{Qwen} ecosystem, deployed on \textbf{ADL Foundry}, the team's Apsara data-lake analytics platform.}
        \resumeItem{Built internal-platform integrations coupling \textbf{Spark}, \textbf{Flink}, and \textbf{Apache Iceberg} lakehouse engines with Data-Juicer Agents, driving \textbf{batch} and \textbf{Kafka}-backed \textbf{streaming} pipelines that turn raw multimodal data into analysis-ready datasets on \textbf{OSS} object storage.}
        \resumeItem{Built an agent-to-\textbf{Ray} distributed execution pipeline over private-cloud object storage (OSS) and \textbf{Kubernetes}, with hands-on \textbf{Kubernetes Operator} development for job lifecycle automation --- \textbf{2,793 records, 16 operators, $\sim$5 minutes} --- landed upstream through a merged PR.}
        \resumeItem{Analyzed \textbf{agent execution traces} to localize context bloat; designed a three-layer context-control architecture (tool-output compaction, memory compression, hard token budgets) cutting single-tool context payloads by \textbf{82\%} and stabilizing multi-turn agents on 30k-token models.}
        \resumeItem{Shipped a full-stack agent workspace (\textbf{React / FastAPI / MySQL}) with persistent sessions, execution tracing, conversation rollback, and multi-instance isolation (\textbf{4 container-image releases}); contributed \href{https://github.com/datajuicer/data-juicer-agents/pulls?q=is\string:pr+author\string:CodyQin}{\textbf{5 merged PRs}} (unit-tested, CI green) to the 7k$\star$ Data-Juicer ecosystem, plus 1 proposal adopted by maintainers.}
      \resumeItemListEnd

    \resumeSubheading
      {Blueberry AI}{Burlingame, CA (Remote)}
      {AI Engineer Intern}{Nov. 2025 -- Jan. 2026}
      \resumeParagraph{Built a \textbf{Milvus + MySQL} feature-extraction and auto-tagging pipeline batch-indexing \textbf{200K+ images} for landmark / IP / SKU retrieval; benchmarked \textbf{DINOv3, Qwen-VL, and BGE-VL} under a unified recall / precision / latency framework; delivered a YOLO + vector-search retrieval POC validated for commercial use.}

    \resumeSubheading
      {Tsinghua Sichuan Energy Internet Research Institute}{Chengdu, China}
      {Algorithm Engineer Intern}{Sep. 2025 -- Oct. 2025}
      \resumeParagraph{Engineered a \textbf{DeepLabV3+ (PyTorch)} segmentation model for microgrid fault diagnosis reaching \textbf{92\% accuracy}; owned the full data lifecycle of \textbf{500+ images} from preprocessing to deployment, anchoring winning infrastructure bids.}

    \resumeSubheading
      {Leiting Games (G-bits)}{Shenzhen, China}
      {AI R\&D Intern (AIGC \& Automation)}{May. 2025 -- Aug. 2025}
      \resumeParagraph{Built a \textbf{Dify + RPA} automated monitoring pipeline (\textbf{98.85\% accuracy}, 1,000+ violations detected) and fine-tuned \textbf{SDXL / Flux LoRA} models driving \textbf{117K+ installs, ¥24.6M billings (196\% ROI)}; led a 5-person team to \textbf{3$\times$} creative output (2025 Best Practice).}

  \resumeSubHeadingListEnd

%-----------PROJECTS-----------
\section{Projects \& Publications}
  \resumeSubHeadingListStart

    \resumeProjectHeading
      {\textbf{ChatDKU --- RAG Campus Q\&A Assistant} $|$ \emph{Python, LlamaIndex, LangChain, Redis}}{Jul. 2024 -- May. 2025}
      \resumeParagraph{Co-founded a university-adopted \textbf{RAG} Q\&A agent indexing \textbf{12,000+ webpages} over a privately deployed LLM; implemented \textbf{hybrid search (BM25 + vectors)}, \textbf{HyDE}, multi-level \textbf{Redis} caching, and \textbf{SSO} web deployment; officially deployed at DKU and documented as a DKU Signature Work.}

  \resumeSubHeadingListEnd

  \vspace{2pt}
  \begin{itemize}[leftmargin=0.15in, label={}, itemsep=0pt, topsep=1pt, parsep=0pt]
    \small{\item{
     \textbf{Publication}: \emph{Privacy-Preserving Room Occupancy Estimation Using Federated Analytics of BLE Packets} --- \textbf{ACM SenSys 2024} (Demo Track, second author), pp. 889--890; \textbf{95\%} real-world accuracy via federated learning over BLE.\vspace{-3pt}
    }}
 \end{itemize}

\end{document}
